\documentclass[12pt]{article}
\usepackage[margin=1in]{geometry}
\usepackage{times}
\usepackage{setspace}
\usepackage{enumitem}
\onehalfspacing
\setlength{\parindent}{0pt}
\setlength{\parskip}{4pt}

\begin{document}

\begin{center}
\textbf{\large Individual Contribution Report --- Project 2}\\[4pt]
TransferDB Application --- CMPE 321, Spring 2026
\end{center}

\vspace{2pt}

\textbf{Personal Information.}
Name: Kemal \"{O}nal \quad Student ID: 2021400018 \quad Group Number: 01

\textbf{Tasks \& Contributions.}
I owned the analytical, documentation, and quality-assurance side of
Project 2. My primary deliverable was \texttt{part3.pdf}, the Schema
Refinement and Normalization report: I enumerated the non-trivial
functional dependencies for every relation strictly from the rules
stated in the project specification (following forum \#2's guidance
that FDs must come from explicit constraints rather than real-world
assumptions), including the family of candidate keys
\texttt{(stadium\_ID, match\_datetime)},
\texttt{(referee\_ID, match\_datetime)}, and the two club variants
that the 120-minute rule forces into existence. I then analyzed each
relation for BCNF compliance and identified
\texttt{stadium\_ID} $\to$ \texttt{city} in Club as the single
BCNF-violating dependency, wrote up the tradeoff between decomposing
Club into a BCNF-compliant pair versus keeping the redundancy in
check via triggers, and justified our decision to retain Club in 3NF
so that the C8 ``club city equals stadium city'' rule can be enforced
locally without subquery checks in every insert.

I also compiled the constraints-added catalog (the full mapping from
each business rule to its enforcement mechanism --- CHECK vs.\
trigger vs.\ stored procedure), reviewed the seven triggers and two
stored procedures against the spec to confirm every rule in \S2.2,
\S2.3, and \S4 has a database-level enforcement path, and documented
how the four forum clarifications (strict enum for competition type,
FD listing scope, no suspension carry-over across seasons, and
concurrent red + 5-yellow resolving to a single-match ban with a
yellow-count reset) are reflected in the schema and procedures. I
tracked those clarifications in \texttt{forum\_clarifications.md} so
they stay discoverable for future iterations.

On the QA side I designed the spec-compliance test scenarios --- past
match scheduling, 120-minute conflicts on shared stadium/referee/club,
squad-size boundaries (below 11, exactly 11, above 23), referee-only
result submission with and without the session variables set,
suspended-player squad insertion, and Permanent auto-termination via
the transfer procedure --- and validated each against a running MySQL
container from a direct client. I reviewed the Flask templates for
role-appropriate data exposure and consistency of Turkish labeling,
and I prepared the final submission packaging, ensuring the ZIP
layout matches spec \S6.

\textbf{Collaboration \& Challenges.}
My teammate and I split responsibilities along a clean seam: he
owned schema, triggers, stored procedures, backend, and templates;
I owned the normalization analysis, the constraint catalog, the test
scenarios, and the submission report. We synced via the shared
repository whenever a schema decision affected the write-up (for
example, when the suspension table design changed we reworked the
Match\_Stats BCNF argument to reflect the new FDs). The hardest part
was being rigorous about the FD derivation: real-world football
intuitions kept suggesting FDs that were not actually justified by
spec text, and the forum \#2 guidance forced me to drop several of
them.

\textbf{Self-Assessment.}
Writing Part 3 sharpened my understanding of the difference between
3NF and BCNF in practice, and of why the ``right'' normal form is
sometimes not BCNF when enforcement cost is considered. Working
alongside the implementation also gave me a clearer picture of which
constraints realistically require triggers rather than CHECK.

\textbf{Use of AI Tools.}
No AI tools were used for this project. All code, SQL, and
documentation were written by the team.

\end{document}
